% ============================================================================
% Chapter 10 — Native C++ Scripting
% ============================================================================
\chapter{Native C++ Scripting}

\textbf{Header:} \texttt{script/CppScript.hpp}\\
\textbf{Namespace:} \texttt{Caffeine::Script}

CppScript is the engine's native scripting interface. You write a C++ class
that inherits from \texttt{CppScript}, override the callbacks you need, and
register it with a macro. The engine's \texttt{ScriptSystem} attaches script
instances to entities and calls the lifecycle methods every frame.

\section{CppScript Base Class}

\begin{lstlisting}
class CppScript {
public:
    virtual ~CppScript() = default;

    virtual void onCreate(ECS::Entity entity, ECS::World& world);
    virtual void onUpdate(ECS::Entity entity, ECS::World& world, f32 dt);
    virtual void onDestroy(ECS::Entity entity, ECS::World& world);
    virtual void onCollision(ECS::Entity entity, ECS::Entity other,
                             ECS::World& world);
};
\end{lstlisting}

All callbacks have default empty implementations; override only what you
need.

\begin{longtable}{lp{9.5cm}}
\toprule
\textbf{Callback} & \textbf{When called} \\
\midrule
\texttt{onCreate}    & Once, when the script is attached to an entity. Use to initialise state. \\
\texttt{onUpdate}    & Every frame. \texttt{dt} is the elapsed time in seconds. \\
\texttt{onDestroy}   & Once, when the entity is destroyed or the script is removed. \\
\texttt{onCollision} & When the entity's collider touches another entity's collider. \texttt{other} is the other entity. \\
\bottomrule
\end{longtable}

\section{Writing a Script}

\begin{lstlisting}
#include "script/CppScript.hpp"
#include "ecs/World.hpp"
#include "ecs/Components.hpp"
#include "physics/PhysicsComponents2D.hpp"

class PlayerController : public Caffeine::Script::CppScript {
public:
    void onCreate(ECS::Entity entity, ECS::World& world) override {
        // Nothing to initialise yet
    }

    void onUpdate(ECS::Entity entity, ECS::World& world, f32 dt) override {
        auto* vel = world.get<Velocity2D>(entity);
        if (!vel) return;

        // Move horizontally based on input
        // (see Chapter 9 for InputManager usage)
        vel->x = m_moveDir * m_speed;

        // Simple gravity
        vel->y += 980.0f * dt;  // pixels/s^2
    }

    void onCollision(ECS::Entity entity, ECS::Entity other,
                     ECS::World& world) override {
        // Reduce health if colliding with an enemy
        if (world.has<EnemyTag>(other)) {
            if (auto* hp = world.get<Health>(entity)) {
                hp->current = (hp->current > 10) ? hp->current - 10 : 0;
            }
        }
    }

private:
    f32 m_speed   = 200.0f;
    f32 m_moveDir = 0.0f;
};

// Register the script — place this in the .cpp file, outside any function
REGISTER_CPP_SCRIPT(PlayerController)
\end{lstlisting}

\section{The REGISTER\_CPP\_SCRIPT Macro}

\begin{lstlisting}
REGISTER_CPP_SCRIPT(ClassName)
\end{lstlisting}

This macro registers a factory function in \texttt{CppScriptRegistry} at
static-initialisation time. The editor and scene loader use the registry to
instantiate scripts by name, which is how scripts attached in the editor are
recreated when the game starts.

Place the macro in the \texttt{.cpp} file, at namespace scope, after the
class definition. Do not put it inside a function or header file.

\section{CppScriptRegistry}

\begin{lstlisting}
CppScriptRegistry& reg = CppScriptRegistry::instance();

// Create a script by name (used by the scene loader)
auto script = reg.create("PlayerController");

// List all registered script names (used by the editor)
for (const std::string& name : reg.names()) {
    // display in inspector dropdown
}
\end{lstlisting}

You rarely need to call the registry directly in game code; the
\texttt{ScriptSystem} handles instantiation automatically.

\section{Multiple Scripts on One Entity}

The \texttt{ScriptSystem} supports multiple script instances per entity. To
attach two scripts to the same entity, use the entity inspector in the editor
or add them programmatically:

\begin{lstlisting}
// The ScriptComponent holds a list of CppScript instances.
// Typically the editor handles this, but it can be done in code:
auto& sc = world.add<ScriptComponent>(e);
sc.scripts.push_back(CppScriptRegistry::instance().create("PlayerController"));
sc.scripts.push_back(CppScriptRegistry::instance().create("HealthRegen"));
\end{lstlisting}

\section{Best Practices}

\begin{itemize}[noitemsep]
  \item Keep \texttt{onUpdate} fast. It runs every frame for every entity
        that has this script attached. Expensive work belongs in a background
        job, not in a script callback.
  \item Cache component pointers only within a single callback invocation.
        Component memory can be relocated when other components are added or
        removed from any entity in the world.
  \item Do not destroy the entity inside \texttt{onUpdate}. Schedule the
        destruction for the end of the frame using a deferred command or a
        flag checked by a cleanup system.
  \item Use \texttt{onCreate} for resource setup (loading sounds, caching
        entity references) and \texttt{onDestroy} for cleanup.
\end{itemize}
